\documentclass[11pt]{article}

\usepackage[margin=1in]{geometry}
\usepackage{amsmath,amssymb}
\usepackage{array}
\usepackage{booktabs}
\usepackage{enumitem}
\usepackage{xcolor}
\usepackage{hyperref}
\usepackage{setspace}
\usepackage{float}

\setstretch{1.10}
\setlength{\emergencystretch}{2em}

\hypersetup{
    colorlinks=true,
    linkcolor=blue!50!black,
    urlcolor=blue!50!black,
    citecolor=black
}

% --------------------
% Course information
% --------------------
\newcommand{\course}{DSA4213: Natural Language Processing for Data Science}
\newcommand{\semester}{Semester 1, AY2026/2027}

\begin{document}

\begin{center}
{\LARGE\bfseries Group Project Proposal}\\[0.4em]
{\large Interruptible MLE Agents: Human Intervention in Long-Horizon Tree Search}\\[0.9em]
Group ID: \underline{\hspace{3cm}}\\[0.9em]
\course\\
\semester
\end{center}

% ====================
% Administrative information
% ====================
\section*{Administrative Information}

\begin{tabular}{@{}>{\bfseries\raggedright\arraybackslash}p{0.25\linewidth}
                    >{\raggedright\arraybackslash}p{0.70\linewidth}@{}}
Group members & Nabil Rakaiza Abror (A0000000X) \\
               & Name (Student ID) \\
               & Name (Student ID) \\
               & Name (Student ID) \\[0.7em]
Project direction & Human-in-the-loop, tool-using agents for machine learning
                    engineering. \\[0.7em]
Repository & \url{https://github.com/nabilrakaiza/interruptible-mle-agent} (private during the semester; access granted to the instructor on request) \\
\end{tabular}

% ====================
% Executive summary
% ====================
\section*{Project Summary}

Machine learning engineering (MLE) agents run autonomously for many hours, steering
their search with a validation score they compute themselves. Recent work shows that
this score drifts away from true held-out performance as the run proceeds: selecting
the final submission by test score rather than validation score would raise the medal
rate by 9 to 13 percentage points, and by 16.6 points for greedy agents \cite{aira}.
The agent cannot detect this from inside its own loop, because the signal it would
need is exactly the signal it does not have. A human watching the search could.

We propose an MLE agent built on tree search over candidate solutions that is
\emph{interruptible}: it exposes its search tree, current plan, and distilled
findings in real time, pauses at well-defined safe points, and accepts a small set of
structured human interventions before resuming. The agent is implemented as a small
team of specialised sub-agents with explicit tool permissions, each of which doubles
as an intervention surface. We evaluate on a subset of MLE-bench Lite
\cite{mlebench}, comparing autonomous runs, scripted intervention policies that
simulate a monitoring human, and a small number of live human sessions. We measure
outcome quality, the validation--test gap, interaction cost, and monetary cost. The
project aims to produce the first characterisation of which interventions into an MLE
agent are worth a human's attention, and at what cost.

% ====================
% 1. Problem and motivation
% ====================
\section*{1. Problem, Research Question, and Motivation}

\subsection*{1.1 Research question}

\begin{quote}
\textbf{Which human interventions into a tree-search MLE agent improve its outcomes,
and at what interaction cost?}
\end{quote}

Two sub-questions follow. First, \emph{can a human, given a progress view of the
search, detect and correct the validation--test divergence that the agent cannot
see?} Second, \emph{which kind of intervention matters}: those that change the
agent's next \emph{action} (redirecting the search, pruning a branch, overriding the
final submission), or those that change the agent's \emph{beliefs} (injecting domain
knowledge, correcting a stored conclusion) and leave the decision to the agent?

\subsection*{1.2 Motivation}

The failure we target is measured, large, and currently unaddressed.

\textbf{The loss is quantified.} Toledo et al.\ \cite{aira} tracked what an agent
believes it is achieving (mean 5-fold cross-validation score) against what it
actually achieves on the held-out test set across a 24-hour run. The believed score
keeps climbing while the true score plateaus and eventually declines. Choosing the
final node by test score instead of validation score would raise the medal rate by 9
to 13 absolute points depending on search policy, and oracle final-node selection
alone eliminates the entire gap for agents using Monte Carlo tree search or
evolutionary search. Their mitigation, submitting the top-$k$ nodes ranked by
validation score, recovers roughly 10\% of performance at $k=3$, which is real but is
a blind hedge rather than a judgement.

\textbf{The failure is invisible from inside the loop.} Nam et al.\ \cite{mlestar}
report a case on \texttt{spaceship-titanic} in which validation accuracy rose from
0.819 to 0.868 while test accuracy fell from 0.803 to 0.734, caused by feature
engineering that used the target variable. The agent experienced this as an
improvement. Yamada et al.\ \cite{aisci2} audited their own workshop-accepted paper
after the fact and found roughly 57\% overlap between the training and test sets in
the generated code, along with a figure contradicting the claim it was cited for.
Three human reviewers did not catch either.

\textbf{No existing MLE agent admits a human.} AIDE \cite{aide}, AIRA \cite{aira},
MLE-STAR \cite{mlestar}, ML-Master \cite{mlmaster} and ML-Master 2.0
\cite{mlmaster2} are autonomous by construction: each decides internally when to stop
debugging, which branch to pursue, and what to submit. MLE-STAR notes that a human
can supply a model description or override the target code block, but this is
neither developed nor evaluated. Human-in-the-loop mechanisms for general LLM agents
exist, such as rollback and message editing in AGDebugger \cite{agdebugger} and
co-planning in Magentic-UI \cite{magentic}, but the natural unit of intervention in
an MLE search is a scored node in a solution tree, not a turn in a dialogue.

\subsection*{1.3 Why an agent-based approach}

The task is open-ended code generation from a task description and a raw dataset.
Classical AutoML and neural architecture search require a predefined search space,
which is unavailable here; the LLM-agent framing exists because the search space is
the space of programs \cite{aide}. The system is agentic in the sense studied in this
course: it interleaves reasoning with actions and environment feedback
\cite{react}, searches over candidate solutions with a verifier
\cite{tot}, acts through a constrained agent--computer interface \cite{sweagent},
and maintains reflective memory across a long horizon \cite{genagents}. Our question
is specifically about \emph{agent} autonomy: it arises only because the system takes
hundreds of consequential decisions without supervision.

\subsection*{1.4 What success means}

We will consider the project successful if we deliver:
\begin{enumerate}[leftmargin=*,itemsep=0.15em]
    \item a working interruptible MLE agent that produces valid submissions on the
          majority of our task subset;
    \item a documented intervention vocabulary with defined safe-point semantics,
          organised along the action/belief axis of Section 1.1;
    \item a controlled comparison of autonomous versus intervened runs on outcome,
          process, interaction-cost, and monetary-cost metrics; and
    \item a reproducible artefact: code, configurations, saved search trees, and
          intervention logs sufficient to regenerate every table in the report.
\end{enumerate}
Beating the current state of the art on MLE-bench (56.4\% medal rate
\cite{mlmaster2}) is explicitly \emph{not} a goal; that figure was obtained under
compute budgets far beyond ours. A negative result, in which interventions do not
improve outcomes, is informative and will be reported as such.

% ====================
% 2. Related work and contribution
% ====================
\section*{2. Related Work and Proposed Contribution}

\subsection*{2.1 Benchmark}

MLE-bench \cite{mlebench} comprises 75 offline Kaggle competitions with local grading
code and human baselines from the original leaderboards. The headline metric is the
percentage of competitions in which the agent earns any medal. We use MLE-bench Lite,
the 22 low-complexity competitions, which is the subset most published agent work
reports on and the only tier feasible on student compute.

\subsection*{2.2 MLE agents}

AIDE \cite{aide} reframes MLE as optimisation over the space of Python scripts. Each
node in a search tree is a self-contained script with a scalar fitness, so the prompt
never accumulates the full interaction history. A greedy policy repeatedly applies one
of three operators (draft, debug, improve) to the best working solution.

AIRA-dojo \cite{aira} formalises such agents as a tuple of fitness function,
selection policy, operator set, operator policy, and termination rule, which allows
components to be varied independently. Two findings matter for us. First, the
operator set, not the search policy, is the bottleneck: replacing greedy search with
MCTS or evolutionary search while holding AIDE's operators fixed produced no gain, and
ablating AIDE's memory operator changed almost nothing. Second, agents systematically
overfit their proxy metric (Section 1.2).

MLE-STAR \cite{mlestar} retrieves candidate models from the web to form an initial
solution, then uses generated ablation studies to identify the highest-impact code
block and refines only that block. It also introduces explicit data-leakage and
data-usage checkers, and reaches 43.9\% on MLE-bench Lite with a small, inexpensive
model.

ML-Master \cite{mlmaster} runs MCTS with parallel branches and supplies the reasoning
model only with its parent and sibling nodes. ML-Master 2.0 \cite{mlmaster2}
generalises this into Hierarchical Cognitive Caching, promoting raw traces into
distilled judgements at phase boundaries and into transferable wisdom at task
completion. It is the current state of the art at 56.4\%. Its distilled-judgement
tier is close in form to the reflections of Generative Agents \cite{genagents}, and
we adopt a similar format for our progress view.

FML-bench \cite{fmlbench} argues that these benchmarks measure only final performance
and proposes an exploration-diversity metric over an agent's proposals.

\subsection*{2.3 Human-in-the-loop agents}

AGDebugger \cite{agdebugger} shows that rollback and message editing mitigate
cascading errors in long-horizon multi-agent systems. Magentic-UI \cite{magentic}
supports co-planning, action approval, and answer verification for general digital
tasks. Neither addresses a search over machine learning solutions, where interruption
must respect long-running training jobs and where the artefact under review is a
scored program.

\subsection*{2.4 Proposed contribution}

The tree representation and the draft/debug/improve operators are adapted from AIDE
\cite{aide} with the prompt-level improvements reported by AIRA-dojo \cite{aira}; the
progress summary format is adapted from ML-Master 2.0 \cite{mlmaster2}; the leakage
checker follows MLE-STAR \cite{mlestar}. Our contribution is:

\begin{enumerate}[leftmargin=*,itemsep=0.2em]
    \item \textbf{Interruptibility as a property of the search loop.} A checkpoint
          and resume mechanism with defined safe points, a structured intervention
          vocabulary, and a real-time progress view. No published MLE agent provides
          this.
    \item \textbf{The first measurement of intervention value in MLE.} A controlled
          comparison of autonomous runs, scripted intervention policies, and live
          human sessions, reported against outcome metrics, the validation--test gap,
          and interaction cost.
    \item \textbf{Structural separation of evaluation from optimisation.} The
          evaluation harness is built once by a dedicated sub-agent and is immutable
          to the agent that writes solution code, enforced by file permissions
          rather than by instruction. Prior systems check for leakage after the fact
          \cite{mlestar}; we make the contract structural, stress-test the checker,
          and ablate it.
\end{enumerate}

% ====================
% 3. Agent setup
% ====================
\section*{3. Agent and Task Setup}

\subsection*{3.1 Task and environment}

Given a competition description and a dataset directory, the system must produce
\texttt{submission.csv}. Submissions are graded locally by MLE-bench's grading code
against the competition's medal thresholds. Each run executes inside an isolated
container with a fixed CPU, memory, and GPU quota, a per-run wall-clock budget, and a
per-script execution timeout.

\subsection*{3.2 Sub-agents}

\begin{itemize}[leftmargin=*,itemsep=0.25em]
    \item \textbf{A1, Evaluation agent.} Runs once at the start. Inspects the data,
          constructs three disjoint splits (training, validation, and a small
          \emph{probe} split), and writes the scoring function to a harness module.
          The validation split drives node selection. The probe split is hidden from
          A3 and from selection and is shown only in the progress view, where it gives
          the monitor a signal the agent never optimises against. A1 also screens
          each candidate script for target leakage and unused data sources before
          execution, following \cite{mlestar}.
    \item \textbf{A2, Retrieval agent.} On draft operations, queries the web for
          candidate model families and example usage for the task type, and returns
          them as model cards in a fixed schema.
    \item \textbf{A3, Coding agent.} Implements the draft, debug, and improve
          operators. Produces a short plan plus a self-contained script that trains,
          scores itself through the harness, and writes a submission.
    \item \textbf{A4, Reflection agent.} At phase boundaries, compresses the traces
          of completed nodes into short findings, each with the evidence that
          produced it (node identifiers and scores), and removes the raw traces from
          working context.
    \item \textbf{A5, Tree controller.} Maintains the search graph, selects the next
          node and operator, enforces stopping rules, and owns checkpointing and
          intervention application. It is the only writer to the tree.
\end{itemize}

\subsection*{3.3 Tool permissions}

Permissions are enforced by the environment, not by prompts. Table~\ref{tab:perms}
lists what each component may do.

\begin{table}[H]
\centering
\small
\begin{tabular}{@{}lccccc@{}}
\toprule
& \textbf{Read data} & \textbf{Write harness} & \textbf{Execute code}
& \textbf{Web access} & \textbf{Write tree / KB} \\
\midrule
A1 Evaluation  & yes & yes (once) & yes, sandbox & no  & no \\
A2 Retrieval   & no  & no         & no           & allowlisted & model cards only \\
A3 Coding      & yes & \textbf{no}& yes, sandbox & no  & no \\
A4 Reflection  & no  & no         & no           & no  & findings only \\
A5 Controller  & no  & no         & no           & no  & yes \\
Human monitor  & view & no        & no           & no  & via queue only \\
\bottomrule
\end{tabular}
\caption{Permission matrix. A3 cannot modify the harness or reach the network, so it
can neither redefine its objective nor fetch arbitrary code.}
\label{tab:perms}
\end{table}

\subsection*{3.4 Knowledge base}

A structured store, not a vector index in the first version. It holds model cards
from A2 and findings from A4, each with its source or evidence. AIRA-dojo
\cite{aira} found that removing AIDE's memory operator barely changed performance,
while ML-Master 2.0 \cite{mlmaster2} found its knowledge tier necessary for top
results. The contradiction is unresolved, so we treat memory as a hypothesis to test.
Retrieval over the store will be introduced only if injected context is observed to
saturate, since within a single run the number of findings is small.

\subsection*{3.5 Interruption semantics}

\textbf{Safe points} are node boundaries: after a script has executed and its result
has been scored and written to the tree, and before A5 selects the next node. This is
the only moment at which all state is consistent. At a safe point A5 serialises the
tree, knowledge base, harness, elapsed budget, and random state, then drains an
intervention queue. The human interface is a Streamlit dashboard running as a
separate process from the agent: it reads the checkpoint and appends to the queue,
and never touches live state. The agent is therefore the only writer to the tree,
and the dashboard can refresh, crash, or be opened by several users without
affecting the run. Scripted intervention policies (Section 4.2) use the same queue,
so human and scripted interventions pass through an identical code path. Every
checkpoint is retained, which allows a finished run to be replayed step by step for
analysis and demonstration.

\textbf{Intervention vocabulary}, grouped by what it changes:
\begin{itemize}[leftmargin=*,itemsep=0.15em]
    \item \emph{Action interventions.} \textbf{Prune} a node or subtree;
          \textbf{Redirect} the next node and operator; \textbf{Override} the final
          submission.
    \item \emph{Belief interventions.} \textbf{Inject} a model card or approach into
          the knowledge base; \textbf{Correct} or delete a finding from A4.
    \item \emph{Budget intervention.} \textbf{Rebudget} the remaining wall clock.
\end{itemize}

A running script is not interrupted by default. As an optional budget-saving action,
the monitor may terminate a running script; the node is then marked
\emph{abandoned} and contributes no score, since a partial score is not comparable to
a complete one and would corrupt selection. Wall clock spent on an abandoned node is
counted as part of that intervention's cost. The budget clock does not pause while
the agent waits on a human, so intervened and autonomous runs receive equal time.

Queued actions are applied in timestamp order. An action whose target no longer
exists or is no longer in the expected state is dropped and the drop is logged.
Every applied or dropped action is recorded with its timestamp, target, and optional
note, which yields a complete audit trail of human influence on each run.

\subsection*{3.6 Constraints and stopping conditions}

A2 may fetch model documentation, library examples, and research papers from an
allowlist, and is prohibited from retrieving competition solutions, winning
notebooks, or leaderboard discussion, in line with MLE-bench's rule against looking up
solutions \cite{mlebench}. Search terminates on wall-clock exhaustion or a node cap.
A debug chain terminates after a fixed number of consecutive failures, and an improve
chain terminates after a fixed number of attempts that fail to exceed a small
relative-improvement threshold, following \cite{mlmaster}. Failed nodes remain in the
tree, marked buggy, and are visible in the progress view.

% ====================
% 4. System or study design
% ====================
\section*{4. Proposed System or Study Design}

\subsection*{4.1 Main loop}

\begin{enumerate}[leftmargin=*,itemsep=0.1em]
    \item A1 builds the train, validation, and probe splits and the harness, then
          freezes the harness.
    \item A5 selects a node and an operator (greedy on validation score, with a small
          probability of revisiting a buggy node, following \cite{aide}).
    \item On a draft, A2 supplies model cards; on any operator, the knowledge base
          supplies current findings.
    \item A3 writes a plan and a script. A1 screens it for leakage and unused data.
    \item The script executes in the sandbox and is scored on validation and probe.
          It also writes a test-set submission, which is graded but never shown to
          any agent. The result becomes a new node.
    \item A4 records findings if the node closes a phase.
    \item \textbf{Safe point.} A5 checkpoints and applies queued interventions, then
          returns to step 2.
    \item On termination, the final submission is chosen by the configured policy:
          best validation, top-$k$, or human override.
\end{enumerate}

The part owned by the group is the checkpoint, queue, and progress-view layer, the
evaluation contract between A1 and A3, and the study design. The final report will
include a system diagram and annotated example trajectories.

\subsection*{4.2 Comparisons and ablations}

Table~\ref{tab:conditions} lists the conditions. C2 and C3 require \emph{no
additional runs}: every node retains its script, validation score, and graded test
score, so both are computed post hoc from the saved trees of C1. They bound how much
any selection-time intervention could be worth and reproduce the analysis of
\cite{aira} on our own system.

Live human sessions (C6) are not used as replayable data. LLM outputs are
nondeterministic, so a recorded action such as ``prune node 17'' refers to a
different script on any rerun. Instead, we observe what participants do and why,
express the observed behaviour as explicit rules, and evaluate those rules as
scripted conditions (C4, C5) over fresh seeds. The live sessions therefore serve as a
source of hypotheses and qualitative evidence.

\begin{table}[H]
\centering
\small
\begin{tabular}{@{}llp{0.47\linewidth}l@{}}
\toprule
\textbf{ID} & \textbf{Condition} & \textbf{Purpose} & \textbf{Runs} \\
\midrule
C1 & Autonomous baseline & Greedy search, best-validation submission. & yes \\
C2 & Oracle final-node selection & Upper bound on selection-time gain. & post hoc \\
C3 & Top-$k$ submission ($k=3$) & Blind-hedge comparison from \cite{aira}. & post hoc \\
C4 & Scripted: divergence pruning & Prune a subtree when validation gains are not
      matched on the probe split. Action intervention. & yes \\
C5 & Scripted: finding correction & Delete any finding whose supporting nodes are
      later contradicted. Belief intervention. & yes \\
C6 & Live human sessions & Qualitative runs with think-aloud; source of C4/C5 rules.
      & few \\
\midrule
A1-off & No evaluation agent & A3 writes its own split and metric. Tests the
      structural separation claim. & yes \\
A4-off & No reflection agent & Tests whether accumulated findings help at all. & yes \\
A2-off & No retrieval agent & Tests web-retrieved models. Optional. & if budget \\
\bottomrule
\end{tabular}
\caption{Experimental conditions. One component changes per row, following the
attribution argument of \cite{aira}. C4 and C5 cover one action and one belief
intervention respectively.}
\label{tab:conditions}
\end{table}

% ====================
% 5. Evaluation
% ====================
\section*{5. Preliminary Evaluation Plan}

\subsection*{5.1 Metrics}

\textbf{Outcome}, reported as in MLE-bench \cite{mlebench}: valid submission rate,
above-median rate, any-medal rate.

\textbf{Process}, where our contribution is visible:
\begin{itemize}[leftmargin=*,itemsep=0.15em]
    \item \textbf{Validation--test gap}: medal rate when the final node is selected on
          test score minus medal rate when selected on validation score, per
          condition. This is the quantity the project exists to move.
    \item \textbf{Gap trajectory}: the same quantity over wall-clock time, testing
          whether divergence widens as reported in \cite{aira}.
    \item \textbf{Probe fidelity}: correlation between probe and test scores across
          nodes, which determines whether the monitor's signal is trustworthy.
    \item \textbf{Exploration diversity} \cite{fmlbench}, to test whether
          intervention broadens or narrows the search.
\end{itemize}

\textbf{Interaction cost}: number of interventions per run, their timing, human time
spent per intervention in C6, and wall clock lost to abandoned nodes.

\textbf{Monetary and system cost}: input and output tokens per node and per run, API
cost, wall clock, and GPU-hours, logged for every condition.

\subsection*{5.2 Reliability checks on LLM components}

A1's leakage screen and A4's findings are both LLM judgements, which can be
unreliable \cite{llmjudge}. We therefore run two targeted checks:
\begin{itemize}[leftmargin=*,itemsep=0.15em]
    \item \textbf{Leakage stress test.} We construct a set of scripts with planted
          leakage of known types (target used as a feature, test statistics used in
          imputation, validation rows in training) alongside clean controls, and
          report A1's detection rate and false-positive rate.
    \item \textbf{Finding accuracy.} For a sample of A4 findings, we check whether
          the claim is supported by the cited nodes' scores, and report the fraction
          that is unsupported.
\end{itemize}

\subsection*{5.3 Protocol}

\textbf{Task subset.} Six MLE-bench Lite competitions chosen for small dataset size
and modality coverage: \texttt{spooky-author-identification} and
\texttt{detecting-insults-in-social-commentary} (text), \texttt{random-acts-of-pizza}
and \texttt{nomad2018-predict-transparent-conductors} (tabular),
\texttt{aerial-cactus-identification} and \texttt{leaf-classification} (image). All
are under 40\,MB.

\textbf{Budget.} A 4-hour wall-clock budget per run instead of MLE-bench's 24 hours,
trading absolute performance for seed count, following the argument in \cite{aira}
that a smaller task set with more seeds yields more reliable conclusions.

\textbf{Seeds.} Five seeds per task for C1, C4, C5; three for the ablations.
AIRA-dojo \cite{aira} shows that agent rankings can flip under three and even five
seeds and recommends at least ten. We cannot afford ten, so we report bootstrapped
95\% confidence intervals and do not claim rankings whose intervals overlap.

\textbf{Reproducibility.} LLM outputs are not deterministic even at temperature zero,
owing to provider-side batching and non-associative floating-point reduction, and
generated training scripts carry their own nondeterminism. We therefore make
distributional rather than exact-replay claims. We pin model version identifiers,
fix and record seeds for every component we control, and log every prompt, raw
completion, model version, sampling parameter, and timestamp.

\textbf{Error analysis.} Failed and poor nodes will be coded into a taxonomy (runtime
error, timeout, invalid submission format, leakage, overfitting to validation,
unused data, stagnation), reported per condition. We will include at least two
annotated trajectories, one where intervention helped and one where it did not.

\subsection*{5.4 Mapping claims to evidence}

\begin{table}[H]
\centering
\small
\begin{tabular}{@{}p{0.42\linewidth}p{0.52\linewidth}@{}}
\toprule
\textbf{Claim} & \textbf{Evidence} \\
\midrule
The agent overfits its own proxy metric & Gap trajectory over C1 runs \\
The probe gives the monitor a usable signal & Probe fidelity across C1 nodes \\
A human-style policy can correct the drift & Gap and medal rate, C4 versus C1,
  bounded above by C2 \\
Judgement beats a blind hedge & C4 versus C3 \\
Action and belief interventions differ in value & C4 versus C5 \\
Structural evaluation separation matters & Leakage incidents and gap, A1-off versus
  C1; stress-test detection rate \\
Accumulated findings help, or do not & Medal rate, A4-off versus C1; finding
  accuracy \\
Interventions are worth their cost & Medal rate change against interaction cost and
  lost wall clock \\
\bottomrule
\end{tabular}
\caption{Each claim is tied to a specific comparison.}
\end{table}

% ====================
% 6. Risks and responsible practice
% ====================
\section*{6. Risks, Limitations, and Responsible Practice}

\textbf{Prompt injection through retrieval.} A2 reads untrusted web content, and
text derived from it reaches A3, whose output is \emph{executed}. A malicious or
merely misleading page could therefore steer generated code. Mitigations follow the
instruction-hierarchy principle \cite{insthier} that retrieved content is data, not
instructions: A2 returns only a fixed model-card schema rather than free text; model
cards are delimited and labelled as untrusted in A3's prompt; A2's sources are
allowlisted; and the sandbox has no network access, no credentials, and no access to
the host, so a compromised script is contained. We will red-team this by inserting a
page with an embedded instruction into A2's results and reporting whether it
propagates.

\textbf{Compute.} A full MLE-bench evaluation costs roughly 1{,}800 GPU-hours
\cite{mlebench}. Our plan (Section 7) is far smaller, runs on local machines and
free-tier cloud credits, and all post hoc conditions add no compute. If the budget proves infeasible, we drop tasks before seeds, since seed
count is what protects our conclusions, and drop A2-off first.

\textbf{Reduced budget changes absolute numbers.} AIRA-dojo \cite{aira} observed that
meaningful divergence between agents appears only after roughly 15 hours. Our 4-hour
runs will report lower medal rates than published figures, and our conclusions
concern relative differences between conditions on a shortened horizon.

\textbf{Scripted policies may not represent real users.} C4 and C5 are reproducible
but not behaviourally faithful. C6 mitigates this only partially; with a handful of
participants, drawn from the group, it is illustrative and will not be presented as a
user study.

\textbf{Probe split costs data.} Carving a probe split reduces training data slightly.
It is applied identically in all conditions so comparisons remain fair, but absolute
scores are affected.

\textbf{Reflection may mislead.} A finding written from one noisy run can propagate
silently. Each finding stores its evidence, C5 tests correction, and A4-off bounds
the damage. A null result resolves a genuine contradiction in the literature and will
be reported as a finding.

\textbf{Contamination.} MLE-bench competitions are public and may appear in
pretraining data. The original authors found no correlation between model
familiarity and performance \cite{mlebench}. We additionally restrict A2's sources
and will spot-check final solutions against public discussion as in \cite{mlestar}.

\textbf{Conduct.} Agents execute generated code only inside an isolated container.
API keys are held in environment variables and excluded from the repository. We use
only publicly released MLE-bench data and will not redistribute competition datasets.
Participants in C6 are group members; no external personal data is collected.

% ====================
% 7. Resources and work plan
% ====================
\section*{7. Required Resources and Work Plan}

\subsection*{7.1 Reused versus built}

\begin{table}[H]
\centering
\small
\begin{tabular}{@{}p{0.30\linewidth}p{0.62\linewidth}@{}}
\toprule
\textbf{Reused (documented)} & \textbf{Purpose} \\
\midrule
MLE-bench \cite{mlebench} & Tasks, data preparation, grading code, medal thresholds \\
AIDE open-source code \cite{aide} & Starting point for operators and tree structure \\
Hosted LLM API & All sub-agents; a low-cost model tier, which MLE-STAR showed is
  viable \cite{mlestar} \\
Container runtime & Sandboxed execution \\
Streamlit & Progress dashboard and intervention interface \\
Python ML stack & PyTorch, scikit-learn, LightGBM, and similar inside the sandbox \\
\midrule
\textbf{Built by the group} & \\
\midrule
\multicolumn{2}{@{}p{0.94\linewidth}@{}}{Evaluation agent and probe split;
permission enforcement; checkpoint, queue, and safe-point mechanism; dashboard views and run replay;
scripted intervention policies; reflection store with evidence; logging and cost
accounting; leakage and prompt-injection stress tests; analysis pipeline.} \\
\bottomrule
\end{tabular}
\end{table}

\subsection*{7.2 Compute and cost}

Run-requiring conditions: C1, C4, C5 at 6 tasks $\times$ 5 seeds (90 runs), and
A1-off, A4-off at 6 tasks $\times$ 3 seeds (36 runs), giving 126 runs of up to 4
hours, or at most about 500 machine-hours. In practice this is lower because
small-data tasks terminate early.

\textbf{Source of compute.} Group members' local machines, supplemented by free-tier
cloud compute (for example Modal). The task subset was chosen so that every task is
feasible on CPU: the text and tabular tasks need no GPU, \texttt{leaf-classification}
ships precomputed tabular features, and \texttt{aerial-cactus-identification} uses
$32\times32$ images small enough for CPU training. A GPU therefore speeds up the two
image tasks but is not required. Runs are independent, so CPU runs are spread across
members' machines and GPU-accelerated runs are sent to cloud credits. Because free
tiers are capped, the plan does not depend on them: if credits run out, the remaining
runs execute on CPU, and if total machine time is still insufficient we apply the
fallback in Section 6.

\textbf{API cost.} The LLM API is expected to cost more than compute. It will be
estimated from measured tokens per node in a pilot run during the first build week,
and reported with per-condition totals in the final report.

\subsection*{7.3 Timeline}

\begin{table}[H]
\centering
\small
\begin{tabular}{@{}p{0.20\linewidth}p{0.72\linewidth}@{}}
\toprule
\textbf{Dates} & \textbf{Milestone} \\
\midrule
29 Sep -- 11 Oct & MLE-bench setup on the six tasks; A3 and A5 baseline; A1 harness
  and probe split; pilot C1 run for cost estimation. (Quiz on 5 Oct.) \\
12 Oct -- 18 Oct & Checkpoint and queue layer; progress view; A4; permissions.
  Lighter week (individual assignment due 16 Oct). \\
19 Oct -- 28 Oct & Full C1, C4, C5 and ablation runs; C6 live sessions; stress tests.
  \\
29 Oct -- 12 Nov & Analysis and presentation. If the group is assigned the earliest
  slot (29 Oct), the presentation will report pilot and partial results and the
  full analysis follows in the report. \\
13 Nov & Final report. \\
\bottomrule
\end{tabular}
\end{table}

\subsection*{7.4 Tentative roles}

\begin{itemize}[leftmargin=*,itemsep=0.1em]
    \item \textbf{Member 1, infrastructure and experiments}: MLE-bench setup,
          sandbox, permissions, compute, and run orchestration.
    \item \textbf{Member 2, core agent}: A3 and A5, operators and search loop.
    \item \textbf{Member 3, evaluation and memory}: A1, probe split, A4, and the
          leakage and prompt-injection stress tests.
    \item \textbf{Member 4, interaction layer}: checkpointing, queue, progress view,
          and the C6 sessions.
\end{itemize}
Statistical analysis, error analysis, and writing are shared. Roles will be assigned
to named members, and a contribution statement will accompany the final report as
required.

% ====================
% Bibliography
% ====================
\begin{thebibliography}{99}

\bibitem{mlebench}
J.~S. Chan et al.
\newblock MLE-bench: Evaluating Machine Learning Agents on Machine Learning
  Engineering.
\newblock \emph{ICLR}, 2025. arXiv:2410.07095.

\bibitem{aide}
Z.~Jiang, D.~Schmidt, D.~Srikanth, D.~Xu, I.~Kaplan, D.~Jacenko, and Y.~Wu.
\newblock AIDE: AI-Driven Exploration in the Space of Code.
\newblock arXiv:2502.13138, 2025.

\bibitem{aisci2}
Y.~Yamada, R.~T. Lange, C.~Lu, S.~Hu, C.~Lu, J.~Foerster, J.~Clune, and D.~Ha.
\newblock The AI Scientist-v2: Workshop-Level Automated Scientific Discovery via
  Agentic Tree Search.
\newblock arXiv:2504.08066, 2025.

\bibitem{aira}
E.~Toledo, K.~Hambardzumyan, M.~Josifoski et al.
\newblock AI Research Agents for Machine Learning: Search, Exploration, and
  Generalization in MLE-bench.
\newblock arXiv:2507.02554, 2025.

\bibitem{mlestar}
J.~Nam, J.~Yoon, J.~Chen, J.~Shin, S.~\"O. Ar{\i}k, and T.~Pfister.
\newblock MLE-STAR: Machine Learning Engineering Agent via Search and Targeted
  Refinement.
\newblock arXiv:2506.15692, 2025.

\bibitem{mlmaster}
Z.~Liu, Y.~Cai, X.~Zhu et al.
\newblock ML-Master: Towards AI-for-AI via Integration of Exploration and Reasoning.
\newblock arXiv:2506.16499, 2025.

\bibitem{mlmaster2}
X.~Zhu, Y.~Cai, Z.~Liu et al.
\newblock Toward Ultra-Long-Horizon Agentic Science: Cognitive Accumulation for
  Machine Learning Engineering.
\newblock arXiv:2601.10402, 2026.

\bibitem{fmlbench}
Q.~Zou, H.~H. Lam, W.~Zhao et al.
\newblock FML-bench: Benchmarking Machine Learning Agents for Scientific Research.
\newblock arXiv:2510.10472, 2025.

\bibitem{agdebugger}
W.~Epperson et al.
\newblock Interactive Debugging and Steering of Multi-Agent AI Systems.
\newblock \emph{CHI}, 2025.

\bibitem{magentic}
H.~Mozannar et al.
\newblock Magentic-UI: Towards Human-in-the-loop Agentic Systems.
\newblock Microsoft Research, 2025.

\bibitem{react}
S.~Yao, J.~Zhao, D.~Yu, N.~Du, I.~Shafran, K.~Narasimhan, and Y.~Cao.
\newblock ReAct: Synergizing Reasoning and Acting in Language Models.
\newblock \emph{ICLR}, 2023.

\bibitem{tot}
S.~Yao, D.~Yu, J.~Zhao, I.~Shafran, T.~L. Griffiths, Y.~Cao, and K.~Narasimhan.
\newblock Tree of Thoughts: Deliberate Problem Solving with Large Language Models.
\newblock \emph{NeurIPS}, 2023.

\bibitem{sweagent}
J.~Yang, C.~E. Jimenez, A.~Wettig, K.~Lieret, S.~Yao, K.~Narasimhan, and O.~Press.
\newblock SWE-agent: Agent-Computer Interfaces Enable Automated Software
  Engineering.
\newblock \emph{NeurIPS}, 2024.

\bibitem{genagents}
J.~S. Park, J.~C. O'Brien, C.~J. Cai, M.~R. Morris, P.~Liang, and M.~S. Bernstein.
\newblock Generative Agents: Interactive Simulacra of Human Behavior.
\newblock \emph{UIST}, 2023.

\bibitem{llmjudge}
L.~Zheng et al.
\newblock Judging LLM-as-a-Judge with MT-Bench and Chatbot Arena.
\newblock \emph{NeurIPS Datasets and Benchmarks}, 2023.

\bibitem{insthier}
E.~Wallace, K.~Xiao, R.~Leike, L.~Weng, J.~Heidecke, and A.~Beutel.
\newblock The Instruction Hierarchy: Training LLMs to Prioritize Privileged
  Instructions.
\newblock arXiv:2404.13208, 2024.

\end{thebibliography}

\end{document}
